\head{第二十版复校：第12节奇异零点与退化层}
\source{法文\texttt{EGA5\_francais.pdf}物理第69--71页；\\英文\texttt{EGAVaiello.pdf}第46--47页；第十九版\texttt{v\_dot\_middle.tex}第5批后段、第6批及第7批前段}

本段从法文第69页的“Proposition 6.1”起，至第71页的“on peut oublier V''”止；紧接的“D'autre part, dans $X_{\mathrm{Grass}^{m}}-V''$”不在本模块内。法文这三页、英文这两页均已实际查看页面图像；又由原始法文PDF渲染了第69页末、第70页末、第71页上部和维数式附近的放大图。Dot第119批及旧中文只作候选译文与历史记录，不充当原页已核验的证据。本报告中的补充证明是本次译校所写，不冒充原稿已有的证明或所引补篇引理。

\begin{longtable}{p{.18\textwidth}p{.17\textwidth}p{.51\textwidth}}
\toprule 法文物理页 & 英文页 & 本模块内容\\\midrule\endhead
69中下部 & 46上半 & 6.1两处维数界；7.1--7.4修改指示；8.2、8.3及8.4--8.6练习；页末手写旁白。\\
70上中部 & 46中部 & 自8.6起推广的写作计划及是否收入《基础》的讨论。\\
70下部至71上部 & 46下部至47上部 & 奇异零点、切空间超额相交、光滑层的闭包、$V^1$与$V''$。\\
71中部 & 47上中部 & 交维数过滤、补篇引理引用、$V^{(k)}$维数式及参数界。\\
71中下部 & 47中部 & $k=2$的像余维估计，止于可不考虑$V''$。\\
\bottomrule
\end{longtable}

\audit{V20-D01}{第6、7节的修改指示}{法英主公式相同；保留原假设与记号作用范围}
法文第69页给出$\dim f(X_i)\geq m+1$和$\dim f(Z)\leq m-1$，英文第46页相同。这是对6.1两处条件的替换，不是一个省去其他假设的新命题。6.1原条目的固有性和连通性条件不能借此消去。7.1仍照用，7.2、7.3作相应修改，7.2中的$m$可能引起混淆，以及7.4证明可直接取余维$m=n$的线性截面，均按原序保留。

法文\emph{laïus}、\emph{mutatis mutandis}及“逐步推进”的写作语气也已译出。英稿对\emph{laïus}所添语种疑问和第69页标记被判为编辑文字，没有误排为数学公式或条目。旧注D17-M025、D17-M026的实质限制保留于正文脚注V20-D01；本次不声称已经验证这些推广后各命题的完整证明。

\audit{V20-D02}{8.2超额维数门槛}{法英同有旧项漏写；新项读法仍记录叠写边界}
法文第69页与英文第46页所写被替换项均为$\dim Y_\xi>\dim X$。英文第21页8.2原条目实际为$\dim Y_\xi>\dim X-1$，故这一处是法稿已经有的漏项，并非英文独有错误。正文不再引用错误旧式，而指称“其中的超额维数条件”。目标公式取$\dim Y_\xi>\dim X-m$，与一般余维的预期维数一致；但法稿减号后的字母有叠写，英编者附\emph{illegible, ask AG}，其读字状态仍在脚注明示。数学上采用正确阈值，不等于声称叠写笔迹已完全辨清。

后面$\dim f(X_i)>0$改为$\dim f(X_i)\geq m$、8.2后续作类似修改、8.3保留$m=1$，均完整保留。该段不用于抹去8.2原条目其他尚待复校的读字或证明问题。

\audit{V20-D03}{页末手写练习旁白}{已恢复可辨法文；区分作者旁白与英译者猜测}
法文第69页手添页码30，与原有31、32组成连续三页。页末手写为\emph{il n'est d'ailleurs sans doute pas indispensable de faire cet exercice, s'il te barbe de le faire}，其意思是“如果做这道练习使你厌烦，想来也并非非做不可”。此前打字句中的\emph{demandant sans doute un peu plus de soin}也已译出，未只保留一个省略的“留作练习”。英稿“Dieudonné or Blass [Tr]”是编者关于受话者的推测，不能倒置为法文作者原注。

\audit{V20-D04}{“自8.6起”与“除8.6外”}{确证英文方向误译；正文按法稿更正}
法文第70页为\emph{tous les développements du n\textsuperscript{o}8 à partir de 8.6}，其后手添\emph{également}。英文第46页的\emph{except 8.6}将方向译反。旧中文D17-M029继承英文，并解释为与上一段的练习范围不同；本次撤销这一归责，正文改成“第8节自8.6起的各项论述，基本上也能推广”。

随后关于冗长的编写工作、尚不知本质依赖一般化的应用、是否收入《基础》、一般框架有助于\emph{dévisser}、几何语境中的形式表述练习、未来也许用到以及决定交由受话者作出的各层意思，均保留原序，没有只摘出数学计划而删除作者的取舍说明。页首被密集划销的旧段不据相邻保留文字逆造，另见V20-D11。

\audit{V20-D05}{奇异零点的概形结构与纤维}{子式与Fitting描述可核验；不得无条件认作每个纤维自身的奇异子概形}
以下是本次补证的精确范围。设$S$为预概形，$P=\mathbf P(E)$，$E$局部自由秩为$r$；记$N=r-1$。设$X\to S$光滑，在所讨论开闭部分相对维数恒为$n$，$f\colon X\to P$为$S$上的无分歧态射，并取$0\leq m\leq n\leq N$。令
\[
B=\operatorname{Grass}^{m}(P),\qquad
Y=X^{(m)},\qquad
\pi\colon Y\to X,\quad p\colon Y\to B.
\]
无分歧性给出满射
\[
f^*\Omega_{P/S}\twoheadrightarrow\Omega_{X/S}.
\]
因目标局部自由，该满射局部分裂。对偶后，$A=T_{X/S}$成为$D=f^*T_{P/S}$中秩为$n$的局部直和因子；$D$的秩为$N$。这说明“$T_x$视为$P$的线性切空间”所需的切空间单射确实成立，而不是由$f$为任意态射得到。

由前段的关联函子及Euler正合列，有
\[
Y\simeq\operatorname{Grass}_{m}(F_X^\vee)
\simeq\operatorname{Grass}_{m}(D),
\qquad F_X^\vee\simeq D\otimes\mathcal O_X(-1).
\]
第二个同构用到格拉斯曼商丛在将原向量丛与可逆模张量时的典范对应。若$G$是$B$上来自$E^\vee$的泛商，则在$Y$上相应泛商是
\[
q\colon\pi^*D\twoheadrightarrow Q,
\qquad Q=(G\otimes\mathcal O_X(1))|_Y.
\]
局部平凡化后，截面$\phi^m$在$X$方向的微分的对偶正是
\[
q|_{\pi^*A}\colon\pi^*A\longrightarrow Q.
\]
可直接在仿射坐标中验证：把线性截面的$m$个方程记为$\ell_i$，$\phi^m$在$Y$上为零，其沿$X/S$的微分就是$(d\ell_i|_{T_{X/S}})_i$；换基或改变线束平凡化时，因已限制在零点概形上，不产生影响该同态的附加项。

为区别于尚未恢复的补篇定义，先把由该矩阵的$(m-k+1)$阶子式定义的闭子概形记为$D_k\subset Y$，规定$D_0=Y$。当$1\leq k\leq m+1$时，其几何点条件是
\[
\operatorname{rank}(q|_A)\leq m-k,
\qquad\text{即}\qquad
\dim(T_x\cap L)\geq n-m+k.
\]
这里的交维数是射影交维数；将过$f(x)$的向量直线除去后，$T_x$给出$D_x$中的$n$维子空间，$L$给出$\ker q$中$N-m$维的方向子空间，二者方向交的维数恰等于原射影交的维数，不应另加或另减1。

相对余切层有局部表示
\[
Q^\vee\longrightarrow\pi^*\Omega_{X/S}
\longrightarrow\Omega_{Y/B}\longrightarrow0.
\]
因此，按Fitting理想的定义，
\[
D_k
=V\!\left(\operatorname{Fitt}_{n-m+k-1}
                 (\Omega_{Y/B})\right)
\quad(1\leq k\leq m+1).
\]
这个等式是由同一表示矩阵的子式得出的概形等式，不是由几何点相同推断。子式理想及余切层均与基变换相容，故对于任何$S'\to S$，该构造的基变换就是$X_{S'}$上的同一构造；对于任意$x\in X$，$D_k\times_X\Spec\kappa(x)$也是在$\operatorname{Grass}_{m}(D_x)$内由同一阶子式定义的Schubert退化概形。

同样，对任意$\xi\in B$，有概形等式
\[
Y_\xi=X_{\kappa(\xi)}\times_{P_{\kappa(\xi)}}L_\xi,
\]
以及
\[
(D_k)_\xi
=V\!\left(\operatorname{Fitt}_{n-m+k-1}
                 (\Omega_{Y_\xi/\kappa(\xi)})\right).
\]
这些公式并不宣称$Y_\xi$总是预期维数的局部完全交。就$D_1$而言，$Y\to S$和$B\to S$都光滑，且$Y$由普遍的$m$个线性方程在$X\times_SB$中切出；相对Jacobian判据说明$p$在一点光滑当且仅当上面的$X$方向微分秩为$m$。于是$|D_1|$确是$p$的非光滑点集，即正文所称截面的奇异零点集。

但若一个特殊截面超额相交，$|(D_1)_\xi|$可能严格大于$Y_\xi$自身在$\kappa(\xi)$上的非光滑点集。例如在域$K$上取$X\simeq\mathbf P^1_K$为$\mathbf P^2_K$中的一条直线，$m=1$，并令$L_\xi=X$。这时$Y_\xi=X$本身光滑，但限制的线性方程恒为零，其微分也恒为零，所以$(D_1)_\xi=X$；等价地，$\operatorname{Fitt}_0(\Omega_{X/K})=0$。在纤维具有预期余维$m$的局部完全交情形，$\operatorname{Fitt}_{n-m}$正好给出该纤维通常的相对非光滑闭子概形。正文因此坚持使用“奇异零点”，未将这两种不同断言混同。

\audit{V20-D06}{经过指定点与光滑层闭包}{法文可辨，撤销旧缺文；保留闭包与光滑性的区别}
法文第70页末手添\emph{passant par x et qui}；第71页首接\emph{est seulement l'adhérence de la sous-variété lisse}。这两处足以恢复完整句法：经过$x$且与$T$超额相交的$L$所组成的簇，是交维数恰比$n-m$大一的光滑子簇的闭包。英文第46页的\emph{only the loneve}不是原法稿缺失的数学概念。

旧注D17-M031对两处难辨及残词所作的保留，现由实见法文撤销；它关于“不凭残句称闭包光滑”的谨慎仍保留。原句中的“簇的闭包”不足以推出任何指定非约化概形的等式。本报告在V20-D08给出的矩阵论证证明精确秩层的光滑性及其在闭退化集合中的稠密性，未借此声称所有可能的厚化都具有同一结构。

\audit{V20-D07}{上标1、撇号及光滑基底}{法稿为$V^1$和$X$；英文为$V'$和$S$；候选错误归责已撤销}
法文第70页先给出$X^{(m)}=V^0$和$X^{(m)\mathrm{sing}}=V^1$；第71页随后两处均是上标$1$：$V^1-V''$和\emph{points de $V^1$ lisses sur X}。第一处的基底也明写$X$。英文第47页排的是$V'$和$S$，而不是$V^0$。这是一项真实的记号重排和一项基底误译，但在$X/S$光滑的假设下，“在$S$上光滑”可从“在$X$上光滑”推出，故不能把英文句子另判为一个数学反例已经否定的命题。

本次正文统一采用这两页法稿的$V^1$，说明英文$V'$是同一对象的别写；$V''$仍是双撇号，$V^{(k)}$仍是括号上标。后续模块的$V^2$若指分歧子预概形，仍不得因形似而与$V^{(2)}$等同。Dot第119批关于“英稿转写为$V^0$”的描述与英文实际图像冲突，已明确撤销；旧注D17-M032对英文撇号的识别原本正确，只是当时尚未见法稿的上标1。

作者对补篇$V''$的集合识别两次使用\emph{sauf erreur}。本次只验证自然行列式结构下的相应几何事实，不据此宣布补篇的超奇异集合、概形结构及所有依赖均已审定。

\audit{V20-D08}{过滤各阶的非光滑点}{从首个退化项起的支集等式可证；完整概形等式不能凭此加强}
继续使用V20-D05的$A\subset D$，并先在一个几何点$x$上工作。设$K=\kappa(\bar x)$，$D_x$为$N$维向量空间，$A_x$为$n$维子空间。记
\[
K_0=\min(m,N-n).
\]
对$1\leq k\leq K_0$，要证明$D_k\to X$的非光滑点支集恰为$|D_{k+1}|$，其中$D_{K_0+1}$为空。此结论依赖子式定义的概形结构；若只给出一个点集而不规定结构，无法谈论其相对光滑性。

固定泛商$q\colon D_x\twoheadrightarrow Q$的一个点，令$\rho=\operatorname{rank}(q|_{A_x})$。先在$A_x$中选择$\rho$个向量，其像构成$q(A_x)$的基；再从$D_x/A_x$方向选择$m-\rho$个向量，其像补成$Q$的基。取这$m$列为枢轴的格拉斯曼仿射图，将它们固定为单位矩阵。延长所选基后，$q|_{A_x}$在整个图上写成
\[
\begin{pmatrix}
I_\rho&U\\
0&W
\end{pmatrix},
\qquad
W\in\operatorname{Mat}_{m-\rho,n-\rho},
\]
其中$W$的各项是相互独立的仿射坐标，在所选点均为零；图中其余自由坐标不参加秩条件。对矩阵作可逆行列运算得到理想等式
\[
I_{m-k+1}(q|_{A_x})
=I_{m-k-\rho+1}(W).
\]
这一步是局部理想的直接化简，给出的正是行列式概形结构。

若$\rho=m-k$，右端由$W$的全部一次坐标生成，故$D_k$在该点附近是仿射图中的坐标子空间，光滑且余维为
\[
(m-\rho)(n-\rho)=k(n-m+k).
\]
若$\rho<m-k$，右端的子式阶数至少为2。其在$W=0$处没有一次项，所以$D_k$的切空间等于该环境仿射图的整个切空间。同时$k\geq1$且$n\geq m$保证所要求的子式确实存在并给出非零方程，因此该闭子概形在这里具有正余维；切空间维数严格大于局部维数，不能在该点光滑。可在同一个图内把$W$变到允许的最大秩$m-k-\rho$，并以标量参数使其趋于零，故精确秩$m-k$的层在闭退化集合中稠密。这也核验了V20-D06的闭包陈述。

上述构造在$X$上局部化后仍适用：$A\subset D$是局部可分裂的子丛，选适配标架即把各$D_k$写成$X$与固定行列式概形的局部乘积。精确秩层在$X$上光滑；更低秩点的几何纤维不光滑，因而$D_k\to X$在那里也不光滑。由此得到的是
\[
\{y\in D_k:\ D_k\to X\text{在}y\text{不光滑}\}
=|D_{k+1}|\qquad(1\leq k\leq K_0).
\]
若$k=0$，左端为空，因为$D_0=Y\to X$是格拉斯曼丛；右端$|D_1|$却可能非空。故原稿笼统的“下一阶”需要正文加入的起始阶限制。当$K_0=0$时，不存在非空的首个退化项，上述断言无须应用。

还须区分这个支集等式与某个指定的奇异子概形等式。以下反例说明不可自动加强。取通用$2\times3$矩阵
\[
M=\begin{pmatrix}a&b&c\\d&e&f\end{pmatrix},\qquad
R=K[a,b,c,d,e,f]/(ae-bd,af-cd,bf-ce).
\]
秩至多1的概形$Z=\Spec R$维数为4，其更深退化层是由$(a,b,c,d,e,f)$定义的约化原点。但是若用$\operatorname{Fitt}_4(\Omega_{Z/K})$定义奇异子概形，得到的理想由三条二次关系的Jacobian矩阵的$2\times2$子式生成。这些生成元全是二次式，且包含$a^2$；而$R$没有一次关系，故该理想不包含$a$，不等于$(a,b,c,d,e,f)$。两者支集均为原点，概形结构却不同。这个矩阵模型可出现在本段的格拉斯曼图中，例如取$m=2,n=3,N=5,\rho=0$。法文实际只说\emph{points non lisses}，本次没有将概形结构相等这一额外断言归给原作者。

英文脚注61、62分别要求询问A.G.及定位补篇引理，已保留在正文。这里给出的证明不充当该引理已经找到的证据。

\audit{V20-D09}{退化层维数式及非空参数}{重新验证旧订正；错误在原适用界，不在维数多项式}
维数计算现在固定在域$K$上的非空光滑纯维$n$纤维，$P$的维数为$N=r-1$，$0\leq m\leq n\leq N$。以下证明重新检查第十九版D17-M-E02的每一步，并沿用其中成立的结论。

在固定点$x$，经过$f(x)$的余维$m$线性子簇$L$，对应于$N$维方向空间$D_x$中的$d=N-m$维子空间$W=\ker q$。令$A_x$为$n$维切方向子空间，并置
\[
t=n-m+k.
\]
精确交维层$\dim(A_x\cap W)=t$可依次构造如下。先选择$A_x$中的$t$维子空间$J$；再选择$D_x/A_x$中的$(d-t)$维子空间$C$；最后，在商掉$J$之后选择把$C$提升到$D_x/J$的图映射，参数为$\operatorname{Hom}(C,A_x/J)$。局部选择一个分裂即可得到这一仿射丛描述，维数不依赖分裂。因此该层维数为
\begin{align*}
&t(n-t)+(d-t)\bigl((N-n)-(d-t)\bigr)+(d-t)(n-t)\\
&\hspace{2em}=t(n-t)+(d-t)m\\
&\hspace{2em}=dm-k(n-m+k).
\end{align*}
这些选择非空，当且仅当
\[
\max(0,n-m)\leq t\leq\min(n,d).
\]
因$n\geq m$，它等价于
\[
0\leq k\leq\min(m,N-n).
\]
也可以从商图直接看出上界：$q(A_x)$缺少的$k$个商方向须由$D_x/A_x$提供，故$k\leq N-n$，并且当然$k\leq m$；反之可选分解商实现这些秩。

闭条件$\dim(A_x\cap W)\geq t$是有限个精确交维层的并。把$k$增加1时，余维$k(n-m+k)$增加
\[
n-m+2k+1>0,
\]
所以最大维数由最小允许交维$t$的一层取得，闭退化层仍具有维数$dm-k(n-m+k)$。该论证不依赖有限域计数、数值代入或若干低维例子。

当$x$变化时，V20-D05所证子丛包含使整个退化族局部成为同一Schubert退化概形与$X$的乘积。因此在这一纯维纤维上，
\[
\dim D_k=n+m(N-m)-k(n-m+k).
\]
另一方面，参数化$P$中所有余维$m$线性子簇的格拉斯曼簇具有维数$m(N+1-m)$。相减得到
\begin{align*}
\dim B-\dim D_k
&=m-n+k(n-m+k)\\
&=(k-1)(n-m)+k^2.
\end{align*}
因此原稿的多项式是正确的；必须订正的是非空范围及没有明确说出的几何纤维、纯维条件。$k=0$也符合该式，此时$D_0=Y$；$k>K_0$时退化层为空，不能继续套用一个有限的维数等式。若相对维数因分支而异，分别使用相应$n$，不能把一个$n$代入整个混合维数纤维。

原稿只有$k\leq m$、$k\leq r-m$。反例$X=P=\mathbf P^2_K$、$r=3$、$n=2$、$m=k=1$同时满足原界，但$T_x=P$，所以$\dim(T_x\cap L)=1<2=n-m+k$，$D_1$为空。这个反例只用来否定原范围；正确的一般范围由上面的子空间构造证明，而不是由反例外推。

\audit{V20-D10}{二阶退化集合的像余维}{有限型条件下的像闭包估计成立；补篇对象的认同仍有条件}
若$K_0\geq2$，代入$k=2$得到
\[
\dim B-\dim D_2=n-m+4\geq4.
\]
若$K_0<2$，则$D_2$为空，不应用该有限维数等式。现另假设所讨论的几何纤维$X$在$K$上有限型。于是$D_2$也为有限型，只有有限个不可约分支；对每一分支，函数域的包含或有限型态射的维数不等式给出其像闭包维数不大于源分支维数。对这些有限个分支取并，得到
\[
\dim\overline{p(D_2)}\leq\dim D_2,
\qquad
\codim\bigl(\overline{p(D_2)},B\bigr)\geq n-m+4.
\]
因此这一像可以放进余维至少4的闭集合，更可以在只保留参数空间余维小于2处的讨论中排除。有限型条件在这里不能单凭“局部有限型”替代：无限个低维像的并可以稠密，源的维数不必控制其整个像闭包的维数。

若$X$在$K$上固有，则$Y\subset X\times_KB$和$D_2\subset Y$均闭，投影到$B$固有，因此像本身闭。原段没有在此重述固有性，正文采用始终有意义的“像的闭包”，未无条件宣称像闭。法文第71页这一句可辨，英稿在Grass后的难辨标记没有对应的尚未读出的数学公式。

最后，从$D_2$转到原作者的$V''$仍需前面的集合识别。正文明确写“若采用上面对$V''$的集合识别”，保留\emph{sauf erreur}的作用范围。本次并未验证补篇超奇异概形的定义，也未将$V''$、$V^{(2)}$与后文表示分歧的$V^2$混为一谈。

\audit{V20-D11}{边注、删除及旧注保全记录}{可辨内容据原页保留；不可辨删文定位而不逆造}
法文第69页：页码30手添于31、32之前；8.2新公式减项有叠写；页末两行手写旁白已译出。被替换旧公式的漏项及新公式字形的不确定性分别记录，不能合称“英文难辨”。

法文第70页：页首有两行以连续叉字符密集划销的旧文字，完整句法无法可靠复原；不将零星可见片段拼成保留段。正文的\emph{également}、关于$X/S$各分支维数至少$m$的边注、$X^{(m)}=V^0$、$X^{(m)\mathrm{sing}}=V^1$，以及页末\emph{passant par x et qui}均保留。$m$后被划去的旧限制由手添\emph{est quelconque}取代；本次不猜测已删除的不等式。

法文第71页：首段公式后的手写勾记不读作对象$V$；$V^1$的上标、$V''$的双撇、\emph{exactement}边添语和两次\emph{sauf erreur}均保留。引入$V^{(k)}$处有一个被划去后重写的同名记号。正式维数等式之前还有一行密集划销的维数草式，可辨开头为$\dim V^{(k)}$，后部字形受删除符覆盖，不能可靠抄出等号及右端；本版保留这一删除记录及原页定位，没有从最终多项式反算一个“原式”。末段手添\emph{de Grass}$^m$已在译文的参数空间指代中体现。

旧注D17-M025--D17-M035的实质内容已逐项处理：D17-M025、026的假设和符号限制继续保留；027的旧公式问题保留并改正来源归属；028的受话者猜测明确属英编者；029的“except 8.6”及所谓作者范围差异撤销；030的相对奇异性限制保留并补充纤维区别；031的可恢复残句已恢复；032的英文撇号判断保留为英文记录，法稿本体改用$V^1$；033的两条英文脚注及补篇引理未定位状态保留，同时补证自然行列式模型；034与D17-M-E02的正确维数计算重新核验并完整保留；035由单纯代入核对补足为有限型像闭包论证。上述恢复不影响本模块精确止点之后的V20后续译校。
